We designed a clean detection experiment to validate KV-cache geometry
as a confabulation signal under stringent controls. The design addresses
four critical flaws identified through adversarial red-teaming of our
prior work (see \cref{sec:redteam}).

\subsection{Design Principles}

\begin{enumerate}[nosep]
  \item \textbf{Single system prompt}: All trials use the identical
    system prompt (``You are a helpful assistant.'') to eliminate
    prompt-fingerprint confounds.
  \item \textbf{Post-hoc behavioral classification}: Instead of
    \emph{instructing} the model to confabulate (which creates an
    encoding-phase fingerprint), we present prompts that naturally
    elicit confabulation and classify behavior from the response.
  \item \textbf{Generation-only features}: KV-cache features are
    extracted from a windowed slice of the generation phase only,
    excluding the encoding (prompt-processing) phase.
  \item \textbf{Ground-truth verification}: Each trial includes a
    factual ground truth, enabling objective behavioral classification.
\end{enumerate}

\subsection{Prompt Design}

We use 200 prompts in two pools:

\begin{itemize}[nosep]
  \item \textbf{Factual} (100): Questions with verifiable answers across
    10 knowledge domains (math, history, science, geography, biology,
    arts, economics, tech, sports, food). Example: ``What is the
    chemical symbol for gold?''
  \item \textbf{Confabulation-inducing} (100): Questions about fabricated
    entities, future events beyond training data, and impossible
    premises across 3 categories. Example: ``What was the final score
    of the 2027 Super Bowl?''
\end{itemize}

All prompts use the same template structure and similar token counts
($\pm 3$ tokens). The confabulation-inducing prompts are designed so
that a well-calibrated model should express uncertainty (hedge), while
an overconfident model may confabulate a specific answer.

\subsection{Behavioral Classification}

Responses are classified into four categories based on content analysis:

\begin{itemize}[nosep]
  \item \textbf{CORRECT}: Factually accurate answer matching ground
    truth.
  \item \textbf{HEDGED}: Model explicitly expresses uncertainty
    (``I don't know,'' ``I'm not sure,'' ``as of my training data'').
  \item \textbf{AMBIGUOUS}: Epistemic caveats without clear uncertainty
    (``I believe,'' ``I could be wrong''). Excluded from primary
    analysis.
  \item \textbf{CONFABULATED}: Confident specific answer that is
    factually wrong (verified against ground truth).
\end{itemize}

Classification uses pattern matching on the response text with
conservative thresholds. The AMBIGUOUS category was added after
red-teaming identified that epistemic hedges contaminated the
CONFABULATED class, diluting the signal.

\subsection{Primary Comparison: HEDGED vs CONFABULATED}

The primary analysis compares HEDGED and CONFABULATED responses to
confabulation-inducing prompts. This is the cleanest possible comparison:

\begin{itemize}[nosep]
  \item Same prompt pool (confabulation-inducing only).
  \item Same system prompt.
  \item Different epistemic behavior (the model chose to hedge or
    confabulate).
  \item No encoding-phase difference (identical prompt structure).
\end{itemize}

The secondary comparison (CORRECT vs CONFABULATED) is reported but
acknowledged as confounded with task difficulty, since factual and
confabulation-inducing prompts differ systematically.

\subsection{Feature Extraction}

For each trial, we extract:

\begin{enumerate}[nosep]
  \item \textbf{Encoding features}: SVD of the key cache after prompt
    processing (before generation begins).
  \item \textbf{Generation-only features}: SVD of a windowed slice of
    the key cache containing only generation-phase tokens. Window size
    is capped at $W = 80$ tokens to control for extreme length
    variation.
  \item \textbf{MP-corrected features}: Signal rank, signal fraction,
    top-SV excess, spectral gap, and norm-per-token, all computed
    relative to the Marchenko-Pastur threshold $\lambda_+$
    (\cref{eq:mp_threshold}).
  \item \textbf{Raw features}: Effective rank, spectral entropy,
    top-SV ratio, norm-per-token, norm variance, angular spread
    (for comparison with FWL-corrected analysis).
\end{enumerate}

\subsection{Classification Pipeline}

Leave-one-out (LOO) cross-validation with logistic regression:

\begin{enumerate}[nosep]
  \item For each held-out trial $i$:
    \begin{enumerate}[nosep]
      \item Fit a \texttt{StandardScaler} on the remaining $n-1$
        training trials.
      \item Transform both training and test features.
      \item For raw features: fit FWL residualization on training data,
        apply to both training and test (FWL-in-LOO).
      \item Fit logistic regression ($C = 1.0$, max\_iter $= 1000$).
      \item Record predicted probability for the held-out trial.
    \end{enumerate}
  \item Compute AUROC from the $n$ predicted probabilities.
\end{enumerate}

The scaler is fitted \emph{inside} the LOO loop on training data only.
Our prior work fitted the scaler on all data before LOO, introducing
a subtle data leak (see \cref{sec:redteam}).

\subsection{Control Battery}

We apply the following controls:

\begin{table}[H]
\centering
\caption{Control battery for the clean detection experiment.}
\label{tab:controls}
\begin{tabularx}{\textwidth}{clX}
\toprule
\# & Control & Implementation \\
\midrule
C1 & FWL residualization & Regress features on token count, use residuals. Applied inside LOO. \\
C2 & Encoding baseline & Report encoding-phase AUROC. Should be $\approx 0.5$ for primary. \\
C3 & Behavioral compliance & Verify behavioral labels match response content. \\
C4 & Domain breakdown & Report behavioral distribution per knowledge domain. \\
C5 & Token-count-only AUROC & Train classifier on token count alone as ceiling baseline. \\
C6 & Permutation testing & Shuffle labels 1000$\times$, compute AUROC distribution. \\
C7 & MP invariance diagnostic & Regress MP features on $\log(n)$ for full-window trials. $R^2 < 0.05$ is GOOD. \\
C8 & Empirical null validation & Column-wise permutation null for $\lambda_+$ agreement with theoretical. \\
C9 & Scaler-in-LOO & StandardScaler fitted inside each LOO fold. \\
C10 & FWL-in-LOO & FWL regression fitted inside each LOO fold. \\
C11 & AMBIGUOUS exclusion & Epistemic hedges excluded from primary comparison. \\
C12 & Belt-and-suspenders & FWL applied to MP features using generation window as confound. \\
C13 & Signed AUROC & Report direction of discrimination, not just magnitude. \\
C14 & Single-feature AUROCs & Report each feature's individual discrimination. \\
\bottomrule
\end{tabularx}
\end{table}

\subsection{Models and Compute}

Three instruction-tuned models, each run on a single NVIDIA RTX 3090
(24GB):

\begin{itemize}[nosep]
  \item \textbf{Qwen2.5-7B-Instruct} (Alibaba): 7.6B parameters,
    28 layers, GQA.
  \item \textbf{Llama-3.1-8B-Instruct} (Meta): 8.0B parameters,
    32 layers, GQA.
  \item \textbf{Mistral-7B-Instruct-v0.3} (Mistral AI): 7.2B
    parameters, 32 layers, GQA.
\end{itemize}

Generation uses greedy decoding (temperature $= 0$) for deterministic
geometry. Each trial takes approximately 5.5 seconds. Checkpoints are
saved every 10 trials for fault tolerance. Trials with fewer than 10
generated tokens are excluded from analysis, as these represent model
refusals or truncated outputs that do not produce sufficient
generation-phase cache for meaningful spectral analysis.

The \texttt{LyraGeometryExtractor} from the Oracle Harness is used for
feature extraction, ensuring consistency between the detection experiment
and the runtime harness.
